% Mirror of content/frontmatter/kurzfassung.tex — keep the content
% identical (distilled from Section 5.1 and the RQ boxes in 4.7).
Cyber Threat Intelligence (CTI) requires analysts to correlate threat
information from heterogeneous sources under time pressure. Large
Language Models (LLMs) can support this work but face three limits in
direct use: outdated parametric knowledge, hallucination, and missing
provenance.

This master's thesis examines how Retrieval-Augmented Generation (RAG)
can improve the quality, reliability, and operational usefulness of
CTI analysis and threat-hunting workflows. An open-source prototype is
designed, implemented, and evaluated that ingests four curated CTI
sources (NVD, CISA KEV, CISA advisories, MISP; 84,982 documents) into
a provenance-preserving index. Hybrid retrieval combines BM25 and
dense vector search with cross-encoder reranking. Answer generation is
layered: verifiable facts are rendered deterministically from passage
metadata, and a locally deployed 8B model contributes only the
narrative, with enforced, validated citations and explicit abstention
without evidence.

The evaluation compares six configurations, from a non-retrieval
baseline to a hosted ablation, on a frozen index state with 28 queries
across four analyst task types and three runs per configuration,
scored with RAGAS, field coverage, and a four-dimension analyst rubric
complemented by scenario analysis. The central finding is that answer
correctness stays at baseline level while traceability rises from 1.23
to 2.56: on this query set, retrieval buys verifiability rather than
raw correctness.

The scientific contribution comprises workflow-level empirical
evidence for RAG in CTI, the finding that generic RAG metrics and
analyst-oriented measures diverge, and a reference model with eight
empirically grounded design principles. The practical contribution is
a reproducible, locally deployable system together with a quantified
local-versus-hosted trade-off.
